\subsection{Why the next direction changes}
\label{sec:pivot}

The tested all-IWSG candidate is rejected for the two matrix-Gaussian scopes.
Its derivative agrees with its declared finite program, yet its model-score
MSE is higher. In the smaller scope, the descriptive error mean and variance
changed from $-1.08143$ and $2.77816$ for the canonical score to $-1.32100$
and $2.71290$ for Phase 4B. In the larger scope they changed from $-0.98767$
and $2.89289$ to $-1.57710$ and $7.66270$. The second scope therefore shows
both more squared error in the mean and more variation across paths. These
across-path decompositions are descriptive; they do not separately estimate
the conditional bias at each fixed dataset.

The result is not adequately explained by saying that the KDM derivative is
different from the atom derivative. Both are candidates for approximating
the exact model score, and a different estimator could have smaller error.
The more useful question is why the new approximation should be closer to
the model. Adding Gaussian observation covariance changes the observation
factor. Adding a second reset changes future particles. Evaluating narrow
kernels with fixed-location importance gradients can amplify noise. None of
these operations, by itself, corrects the original model-score bias.

The earlier restriction to the executed finite derivative made unrelated
score recursions appear inadmissible. That restriction protected the meaning
of the canonical API, but it was too narrow as a research objective. A
Fisher-score estimator can be a valid candidate for the underlying model
without being the derivative of the existing OT-reset scalar. It must be
named and validated accordingly. Likewise, a centered control variate can
reduce variance without removing the bias of its baseline. The distinction
changes what an experiment must measure; it does not relax the mathematics.

One historical cost argument also needs correction. The Section 3.6 memo
grouped PaRIS with inherently quadratic ancestor calculations. The exact
ancestor sum is quadratic, but PaRIS replaces that sum with sampled backward
ancestors and has an expected linear-cost result under stated mixing and
acceptance assumptions \citep[Section 3.1, Algorithm 2; Section 3.2.4,
Theorem 10]{olsson2014paris}. Those assumptions do not hold automatically in
every repository model, particularly when transition densities are singular
or backward acceptance becomes small. Nevertheless, ``always quadratic''
was an incorrect reason to exclude it. The all-components IWSG experiments
already incurred quadratic work, making a measured comparison with
ancestor-averaged methods more relevant than that earlier exclusion suggests.

The next direction therefore assigns different jobs to three components.
An evaluable mixture proposes particles. The original model densities specify
their importance correction. A posterior-score recursion averages over the
ancestor uncertainty that remains. Selective pathwise and importance-weight
gradients are an additional way to reduce variance when mixture expectations
must be differentiated. This combination has not been tested here. It is
preferred as the next hypothesis because every component has a clear role in
the model-score calculation, not because the existing negative results prove
it will perform better.
